\documentclass[10pt,a4paper]{article}
\usepackage[margin=1.55cm]{geometry}\usepackage{amsmath,amssymb,booktabs,enumitem,tcolorbox,hyperref,xcolor,needspace}
\tcbuselibrary{breakable}
\hypersetup{colorlinks=true,urlcolor=blue,linkcolor=blue}
\setlength{\parskip}{5pt}\setlength{\parindent}{0pt}
\newtcolorbox{keybox}[1]{colback=blue!4,colframe=blue!45!black,title={#1},breakable}
\newtcolorbox{trapbox}[1]{colback=orange!6,colframe=orange!65!black,title={#1},breakable}
\begin{document}
\begin{center}{\LARGE\bfseries P2 (WMA12) May 2024 -- Past-paper review}\\[3pt]
{\large A-Level Maths (Edexcel IAL) | Harry Williams | 10 July 2026}\end{center}
\begin{keybox}{How to use this guide}This is a lesson-review guide, not a replacement for the official paper. The official question paper and mark scheme control wording, marks and answers. The lesson evidence controls what was attempted, corrected, deferred or praised. Work through each question with the question paper covered; then use the corresponding multi-step card to reveal a complete route.\end{keybox}
\section*{1. Binomial expansion and targeted coefficients (Q1, 5 marks)}
For $(1+u)^n$, write only as many terms as the task needs: $1+nu+\binom n2u^2+\binom n3u^3+\cdots$. Here $u=-x/6$ and $n=9$, so
\[ (1-x/6)^9=1-\frac32x+x^2-\frac7{18}x^3+\cdots. \]
For the coefficient of $x^3$ in $(10x+3)(1-x/6)^9$, do not multiply every term. Make an exponent-pair list: $10x\cdot x^2$ and $3\cdot(-7x^3/18)$. Therefore the coefficient is $10-7/6=53/6$. In the lesson, the missed contribution was $10x\cdot x^2$; this is an exponent bookkeeping issue, not a new binomial formula issue. A reliable check is to write the wanted degree above each factor before multiplying.\begin{trapbox}{Exam cue}``Coefficient of $x^r$'' means collect every product whose powers add to $r$. It does not mean take only the $x^r$ term from one factor.\end{trapbox}
\section*{2. Arithmetic series: equations before calculator work (Q2, 7 marks)}
Use $u_n=a+(n-1)d$ and $S_n=\frac n2[2a+(n-1)d]$. The official paper says $u_6=2$ and $S_{10}=-80$, hence $a+5d=2$ and $5(2a+9d)=-80$. Solving gives $d=20$, $a=-98$. For part (b), substitute these values into the original sum formula before solving: $\frac n2[-196+20(n-1)]>8000$, i.e. $5n^2-54n-4000>0$. The positive root is about $34.2$, so the smallest integer is $35$. Check $S_{34}\le8000$ but $S_{35}>8000$. The frozen ledger contains older values for this question; they are not used here because the official paper/mark scheme is controlling. The tutor's transferable feedback still applies: clear setup earns method marks and a final substitution check catches a multiplication slip.\begin{trapbox}{Do not confuse}A root $n\approx34.2$ is not itself an admissible number of terms. The requested answer is the smallest integer satisfying the original inequality.\end{trapbox}
\newpage
\section*{3. Logarithms with domain control (Q3, 6 marks)}
The official equation is $2\log_2(2-x)=4+\log_2(x+10)$. First write the domain: $2-x>0$, $x+10>0$, so $-10<x<2$. Combine the left side as $\log_2(2-x)^2$ and rewrite $4$ as $\log_2 16$. Thus $(2-x)^2=16(x+10)$. Solve the resulting quadratic and test every root in the domain. The lesson evidence says the student explained rejected values clearly; retain that habit. A logarithm equation is never finished until candidate values are checked against the arguments and the base condition. For $\log_a(a^6)$, with $a>1$, the answer is $6$: the base and argument match a power.\section*{4. Remainder, factor and discriminant (Q4, 8 marks)}
Given $f(x)=(x-2)(2x^2+5x+k)+21$, dividing by $x-2$ has remainder $f(2)=21$. The tutor explicitly derived the Remainder Theorem: if $f(x)=(x-a)q(x)+r$, set $x=a$ to obtain $f(a)=r$. For the factor $2x-1$, use the root $x=1/2$, not $x=2$: $f(1/2)=0$, giving $k=11$. Expand to $f(x)=2x^3+x^2+x-1$, then divide by $2x-1$ to get $x^2+x+1$. Hence $f(x)=(2x-1)(x^2+x+1)$. Its discriminant is $1-4=-3<0$, so it contributes no real roots; $f(x)=0$ has exactly one real solution, $x=1/2$. The later ledger segment records this completion after an earlier deferred setup.\newpage
\section*{5. Inequalities need sign justification (Q5, 6 marks)}
Expand $(x-y)^3>x^3-y^3$, cancel like terms and factor to get $xy(3y-3x)>0$. Since the question gives $x,y>0$, $xy>0$, so dividing by $xy$ does not reverse the inequality: $3y-3x>0$, hence $y>x$. This sentence about positivity is a mark-winning logical bridge. For a counterexample outside the positive domain, choose values that satisfy the starting inequality while not having $y>x$; verify both the premise and conclusion rather than merely choosing negative numbers.\section*{6. Trapezium rule and integral relationships (Q6, 9 marks)}
The lesson marks Q6 as deferred after time pressure. For the trapezium rule, list the equally spaced ordinates, use $h=0.3$, and apply $\frac h2[y_0+y_n+2(y_1+\cdots+y_{n-1})]$. The official estimate is $2.34$. Then use algebra before re-integrating: $2^x+2x=(2^x-2x)+4x$, so the related estimate adds $\int_2^{3.5}4x\,dx=16.5$, giving $18.84$. Likewise $2^{x+1}-4x=2(2^x-2x)$, giving $4.68$.\newpage
\section*{7. Circles: exact radius then a geometric comparison (Q7, 6 marks)}
Complete the square: $x^2+y^2+8x-10y=29$ becomes $(x+4)^2+(y-5)^2=70$. Thus $C_1$ has centre $(-4,5)$ and exact radius $\sqrt{70}$. $C_2$ has centre $(5,-8)$ and radius $\sqrt{52}$. Their centre distance is $\sqrt{9^2+(-13)^2}=\sqrt{250}$. Since $\sqrt{250}>\sqrt{70}+\sqrt{52}$ (or compare safely using numerical approximations only for the inequality), the circles neither touch nor intersect. The tutor corrected presentation: when ``exact'' is asked, do not replace $\sqrt{70}$ by a decimal final answer.\section*{8. Trigonometry, units and completion (Q8, 12 marks)}
For part (i), set $u=\cos x$ after using $\sin x\tan x=\sin^2x/\cos x$ only with the required nonzero condition, or follow the mark scheme's algebraic route; check the stated interval and give radians to 3 s.f. For the greenhouse model, work in degrees because the model writes $(kt+18)^\circ$. At $t=6$, solve $\sin(6k+18)^\circ=5/6$ using both sine angles; under $0<k<10$, retain the appropriate value. The maximum is $22^\circ\mathrm C$ because sine reaches 1. Set $kt+18=90$ and convert the decimal-hour answer into a clock time to the nearest minute. The lesson correction was variable consistency: do not insert the inverse-sine value in place of the already-calculated $k$.\section*{9. Calculus and equal signed areas (Q9, 8 marks)}
Write $y=8x^{3/2}-2x^{5/2}$ first. Then $dy/dx=12x^{1/2}-5x^{3/2}=x^{1/2}(12-5x)$; the nonzero stationary-point coordinate is $x=12/5$. For equal areas, use a negative sign before the integral below the axis rather than modulus notation. With antiderivative $F(x)=\frac{16}{5}x^{5/2}-\frac47x^{7/2}$, equality causes the common value at $4$ to cancel, leaving $F(k)=0$. Factor $k^{5/2}[16/5-(4/7)k]=0$; reject $k=0$ from the diagram and obtain $k=28/5$.\section*{10. Geometric models and indexing (Q10, 8 marks)}
For dormice, $t=0$ means the initial 2000, so after six years the model is $2000(1.03)^6\approx2390$ to 3 s.f. The lesson focused on avoiding an accidental $t-1$ exponent: write $t=0,1,2$ terms if the wording feels ambiguous. For the vole model $N=ab^t$, use the two supplied data points to find $b$ from a quotient, then $a$ by substitution, respecting 2 significant figures only at the final reporting stage. The ledger records the setup as partial; the multi-step card retains the official full route while labelling that lesson boundary.\newpage
\section*{Worked reasoning and error-proof routines}
\subsection*{A. Q1 coefficient route, written as an examiner can follow}
The word \emph{hence} authorises using part (a), but it does not authorise copying an answer without selecting the relevant powers. Start with the minimal carry-in
\[1-\frac32x+x^2-\frac7{18}x^3+\cdots.\]
Put a small degree label above each available term. The factor $10x$ has degree 1, so it needs the degree-2 term from the expansion. The factor $3$ has degree 0, so it needs degree 3. There is no third possibility because the first bracket has only degrees 0 and 1. This is a fast completeness proof. Now calculate in a single aligned line:
\[10x(x^2)+3\left(-\frac7{18}x^3\right)=\left(10-\frac76\right)x^3=\frac{53}{6}x^3.\]
The coefficient is $53/6$, not the whole term unless the question asks for the term. A useful self-check is dimensional in the algebraic sense: every retained product must visibly contain $x^3$.
\subsection*{B. Q2: a safe simultaneous-equation layout}
The sixth term is not $a+6d$: the first term is $a$, therefore term 6 is $a+5d$. Write both facts before substituting numbers:
\[a+5d=2,\qquad S_{10}=\frac{10}{2}[2a+9d]=-80.\]
The second equation simplifies to $2a+9d=-16$. Doubling the first equation gives $2a+10d=4$; subtraction gives $d=20$, then $a=-98$. Keep the bracket in the sum formula for part (b): $S_n=\frac n2[-196+20n-20]=n(10n-108)$. The inequality is $10n^2-108n>8000$, or $5n^2-54n-4000>0$. Once a positive root has been found, test consecutive integers in the original sum rather than relying on a rounded display. This check is especially valuable under time pressure because it catches an incorrect leading coefficient immediately.
\subsection*{C. Q3: why domain screening must happen twice}
Before manipulating, $2-x>0$ and $x+10>0$, which gives $-10<x<2$. After combining logs, do not forget that exponentiating preserves equality only after the domain has made both logarithms defined. The quadratic may offer values produced by algebra but prohibited by the original arguments. Substitute a retained candidate into both original log arguments: this is more reliable than checking only the rearranged quadratic. The short part (ii) is a notation test: $\log_a(a^6)$ asks for the exponent to which $a$ must be raised to produce $a^6$, so it is $6$.
\subsection*{D. Q4: connect the two theorems rather than memorising two tricks}
Polynomial division has the form $f(x)=(x-a)q(x)+r$. Substitution of $x=a$ removes the quotient term and leaves $f(a)=r$. That one derivation explains both commands. If asked for the remainder on division by $x-2$, use $a=2$ and calculate $f(2)$. If told $2x-1$ is a factor, first solve $2x-1=0$, hence $a=1/2$, then set $f(1/2)=0$. After factorisation, the discriminant is not optional decoration: $b^2-4ac=-3<0$ explicitly justifies why the quadratic supplies no real solutions. State the conclusion in words as well as writing the value $x=1/2$.
\newpage
\subsection*{E. Q5: a proof is different from an algebraic answer}
The supplied positivity condition is a tool, not background text. Expanding gives
\[(x-y)^3=x^3-3x^2y+3xy^2-y^3.\]
After cancelling $x^3-y^3$, rearrange to $xy(3y-3x)>0$. The step that earns the reasoning credit is: ``Since $x>0$ and $y>0$, $xy>0$; dividing by $xy$ preserves the inequality direction.'' Only then may one conclude $3y-3x>0$, hence $y>x$. In the counterexample, a reader must see both that the original premise holds and that $y>x$ fails. This is why trying random negative values without verification is not a method.
\subsection*{F. Q6: using an estimate consistently}
The trapezium rule estimates an integral from a table; do not claim it is an exact antiderivative result. With six ordinates at equal width $h=0.3$, the end ordinates receive weight 1 and interior ordinates weight 2. Preserve the supplied four-decimal table values rather than replacing them inconsistently with calculator values. In the later parts, the point is algebraic reuse: split a new integrand into the one already estimated plus an exactly integrable simple term. When doubling an integrand, double the prior estimate. These transformations are quicker and less error-prone than running a fresh trapezium calculation.
\subsection*{G. Q7: a complete geometry proof has three quantities}
Completing the square identifies $C_1$ without approximation: $(x+4)^2+(y-5)^2=70$. Read centre $(-4,5)$ and radius $\sqrt{70}$ directly. For $C_2$, read centre $(5,-8)$ and radius $\sqrt{52}$. Then calculate the distance between centres directly: $\sqrt{(5+4)^2+(-8-5)^2}=\sqrt{250}$. Compare that with the sum of radii. Finish with the geometric condition: if centre distance is greater than the sum, the circles are separated, so neither touch nor intersect.
\subsection*{H. Q8--Q10: final-answer checklist}
For trigonometry, check angle mode before inverse sine and use every solution allowed by the interval. For the greenhouse, write degrees consistently and attach $^\circ\mathrm C$ to temperature. A decimal time is not a clock time: multiply its fractional hour by 60 and round the minutes. For the equal-area calculus question, integrate the function of $x$, keep the sign of the below-axis region under control, and reject a root inconsistent with the diagram. For the geometric model, explicitly write $P(0)=2000$, $P(1)=2000(1.03)$ and $P(6)=2000(1.03)^6$ before selecting an exponent. This one line prevents the common $t-1$ shift.
\section*{Quick-recall checklist}\begin{itemize}[leftmargin=*]\item Target coefficient: pair exponents whose sum is wanted.\item Remainder / factor: substitute the zero of the divisor.\item Log equation: state domain, solve, then test.\item Inequality: state the sign before dividing.\item Circle proof: compare centre distance with sum of radii.\item Trig modelling: identify degree/radian mode and every required unit/rounding instruction.\item Sequences: decide whether the clock begins at $t=0$ or term 1.\end{itemize}
\end{document}